\subsection{La Asimetría Regulatoria que Condiciona Toda la Propuesta}\label{la-asimetruxeda-regulatoria-que-condiciona-toda-la-propuesta}

Al 20 de septiembre de 2026 el ordenamiento chileno presenta una asimetría que define el diseño de cumplimiento de este proyecto: \textbf{la Ley N° 21.663 está plenamente operativa, mientras que la Ley N° 21.719 está publicada pero no vigente, y la autoridad que ella crea no existe}.

La \textbf{Ley N° 21.719}, que regula la protección y el tratamiento de los datos personales y crea la Agencia de Protección de Datos Personales, fue publicada en el Diario Oficial el \textbf{13 de diciembre de 2024}. Su \textbf{Artículo primero transitorio} difiere la entrada en vigencia «hasta el día primero del mes vigésimo cuarto posterior a su publicación», esto es, el \textbf{1 de diciembre de 2026} ---\textbf{72 días después de la fecha de corte}--- y ordena que entretanto «seguirá operando la antigua ley N° 19.628» . Se corrige aquí un error de atribución frecuente: \textbf{la fecha de vigencia está en el Artículo primero transitorio, no en el tercero}, que fija las reglas de la primera designación del Consejo Directivo .

La \textbf{Agencia de Protección de Datos Personales no está constituida}: el Senado \textbf{rechazó la nómina de consejeros el 19 de mayo de 2026} por no alcanzar los dos tercios, y no consta el ingreso de una nueva . El Ejecutivo ingresó el \textbf{1 de septiembre de 2026} el \textbf{Boletín N° 18.623-07}, que trasladaría la vigencia a 2027 y ampliaría el Consejo de tres a cinco miembros, pero ese proyecto está \textbf{en primer trámite constitucional, sin votación en comisión ni en Sala, y no es ley} .

\begin{quote}
\textbf{Regla de decisión para el equipo AudIT.} Diversos medios y sitios de cumplimiento dan la postergación por consumada, en tiempo pasado. \textbf{Esas afirmaciones son incorrectas a la fecha de corte.} Mientras el proyecto no se apruebe y publique, \textbf{la fecha de vigencia sigue siendo el 1 de diciembre de 2026}. La propuesta económica se dimensiona sobre esa fecha, no sobre la prórroga en trámite.
\end{quote}

\subsection{Ley N° 21.719: Obligaciones Exigibles desde el 1 de Diciembre de 2026}\label{ley-n-21.719-obligaciones-exigibles-desde-el-1-de-diciembre-de-2026}

La ley \textbf{no deroga la Ley N° 19.628}: la sustituye casi íntegramente por reemplazo de su articulado, conservando los \textbf{Artículos 17, 18 y 19} sobre datos de obligaciones económicas y comerciales .

El \textbf{Artículo 12} consagra el consentimiento como regla general. El \textbf{Artículo 13}, verificado contra el articulado, contiene \textbf{cinco letras}: a) obligaciones económicas, financieras, bancarias o comerciales; b) obligación legal; c) celebración o ejecución de un contrato; d) interés legítimo del responsable; e) defensa de un derecho ante los tribunales . Dos hallazgos tienen consecuencia de ingeniería: \textbf{el interés vital no es base autónoma} ---se canaliza por el Artículo 16 bis--- y \textbf{la letra a) no tiene equivalente europeo}.

La frontera responsable/encargado está en el \textbf{Artículo 15 bis}: actuación conforme a instrucciones, prohibición de tratar los datos para objeto distinto del encargo, contrato con \textbf{objeto, duración, finalidad y tipo de datos} y responsabilidad solidaria por daños \textbf{{[}V{]} parcial}. El inciso sobre delegación no pudo verificarse completo: \textbf{no debe afirmarse una prohibición absoluta de subencargo} .

El \textbf{Artículo 2° letra g)} incorpora al catálogo de datos sensibles la \textbf{situación socioeconómica}, \textbf{sin equivalente en el Artículo 9 del RGPD} ; el régimen especial se despliega en los Artículos 16 y siguientes, incluido el \textbf{16 sexies sobre geolocalización}. Los derechos de los Artículos 4 a 11 suman al estándar ARCO el \textbf{bloqueo (8 ter)}, la \textbf{portabilidad (9)} y la \textbf{oposición a decisiones automatizadas (8 bis)}, y el \textbf{Artículo 11} obliga a resolver en \textbf{treinta días corridos}, prorrogables una sola vez por otros treinta ---\textbf{corridos, no hábiles}---.

Dos hallazgos verificados corrigen supuestos difundidos en el mercado de cumplimiento. Primero, el \textbf{Artículo 14 sexies} obliga a reportar brechas \textbf{«por los medios más expeditos posibles y sin dilaciones indebidas»} y \textbf{no fija un plazo de 72 horas} : esa cifra es un traslado indebido del Artículo 33 del RGPD. Segundo, el \textbf{Artículo 27} admite la transferencia internacional por país adecuado, cláusulas contractuales tipo, normas corporativas vinculantes, modelos certificados o excepciones puntuales, pero \textbf{al 20 de septiembre de 2026 no hay ningún país declarado con nivel adecuado}, porque esa declaración compete a una Agencia inexistente .

El régimen sancionatorio (Arts. 34 bis a 37) escala hasta \textbf{5.000, 10.000 y 20.000 UTM} según la gravedad de la infracción, y su desarrollo íntegro consta en el entregable de marco legal chileno. Su relevancia para la propuesta es la magnitud: la exposición patrimonial que justifica el programa de cumplimiento del capítulo 3.
